\section{Categorías con clivaje sobre~$\mathcal{E}$}
\label{VI.7}

\begin{definition}
\label{VI.7.1} Sea $\cal{F}$ una categoría sobre~$\cal{E}$. Se llama
\emph{clivaje}
\index{clivaje|hyperpage}%
de~$\cal{F}$ sobre~$\cal{E}$ a una función que asigna a todo
$f\in\Fl(\cal{E})$ un funtor de imagen inversa por~$f$ en~$\cal{F}$,
denotado por~$f^*$. El clivaje se dice \emph{normalizado}
\index{clivaje normalizado|hyperpage}%
\index{normalizado (clivaje)|hyperpage}%
si $f=\id_S$ implica $f^*=\id_{\cal{F}_S}$. Se llama
\emph{categoría con clivaje} (\resp \emph{categoría con clivaje
normalizado})
\index{categoría con clivaje (\resp con clivaje normalizado)|hyperpage}%
a una categoría $\cal{F}$ sobre~$\cal{E}$ dotada de un clivaje
(\resp de un clivaje normalizado).
\end{definition}

Es evidente que $\cal{F}$ admite un clivaje si y solo si $\cal{F}$ es
prefibrada sobre~$\cal{E}$, y en tal caso $\cal{F}$ admite un clivaje
normalizado. El conjunto de los clivajes sobre~$\cal{F}$ está en
correspondencia biunívoca con el conjunto de los subconjuntos $K$ de
$\Fl(\cal{F})$ que satisfacen las condiciones siguientes:
\begin{enumerate}
\item[a)] Los $\alpha\in K$ son morfismos cartesianos.
\item[b)] Para todo morfismo $f\colon T\to S$ en~$\cal{E}$ y todo
$\xi\in\Ob(\cal{F}_S)$, existe un único $f$-morfismo en~$K$ de destino
$\xi$.
\end{enumerate}

Para que el clivaje definido por~$K$ sea normalizado, es necesario y
suficiente que $K$ satisfaga además la condición
\begin{enumerate}
\item[c)] Los morfismos identidad de~$\cal{F}$ pertenecen a~$K$.
\end{enumerate}

Los
\marginpar{171}
morfismos que son elementos de~$K$ podrán llamarse
\emph{\og morfismos de transporte\fg}
\index{transporte (morfismo de)|hyperpage}%
para el clivaje considerado.

La noción de isomorfismo de categorías con clivaje sobre~$\cal{E}$ es
clara. Más generalmente, pueden definirse los morfismos de
$\cal{E}$-categorías con clivaje como los funtores de
$\cal{E}$-categorías $\cal{F}\to\cal{G}$ que aplican los morfismos de
transporte en morfismos de transporte.
(Estos son, en particular, funtores cartesianos). De este modo, las
categorías con clivaje sobre~$\cal{E}$ son los objetos de una categoría,
la \emph{categoría de las categorías con clivaje sobre}~$\cal{E}$. El
lector explicitará la existencia de productos, ligada al hecho siguiente:
si una categoría sobre~$\cal{E}$ es producto de categorías $\cal{F}_i$
sobre~$\cal{E}$, provista cada una de un clivaje, entonces $\cal{F}$ queda
dotada del clivaje natural correspondiente. Se deja asimismo al lector
explicitar la noción de cambio de base en las categorías con clivaje.

Denotaremos por $\alpha_f(\xi)$ el morfismo canónico
$$
\alpha_f(\xi)\colon f^*(\xi)\to\xi.
$$
Como ya se ha dicho, es funtorial en~$\xi$, es decir, se tiene una
transformación natural
$$
\alpha_f\colon i_Tf^*\to i_S,
$$
donde, para todo $S\in\Ob(\cal{E})$, $i_S$ denota el funtor de inclusión
$$
i_S\colon\cal{F}_S\to\cal{F}
$$

Consideremos ahora morfismos
$$
f\colon T\to S\quad\text{ y }\quad g\colon U\to T
$$
en~$\cal{E}$, y sea $\xi\in\Ob(\cal{F}_S)$. Existe entonces un único
$U$-morfismo
$$
c_{f,g}(\xi)\colon g^*f^*(\xi)\to(fg)^*(\xi)
$$
que hace
\marginpar{172}
conmutativo el diagrama
$$
\xymatrix{
f^*(\xi) \ar[d]_{\alpha_f(\xi)} & & \ar[ll]_{\alpha_g(f^*(\xi))}
g^*(f^*(\xi)) \ar[d]^{c_{f,g}(\xi)} \\
\xi & & \ar[ll]^{\alpha_{fg}(\xi)} (fg)^*(\xi) }
$$
(en virtud de la definición de~$(fg)^*(\xi)$). Para $\xi$ variable,
este morfismo es funtorial; es decir, se tiene una transformación natural
$$
c_{f,g}\colon g^*f^*\to(fg)^*
$$
de funtores $\cal{F}_S\to\cal{F}_U$. Observemos inmediatamente:
\begin{proposition}
\label{VI.7.2} Para que la categoría con clivaje $\cal{F}$
sobre~$\cal{E}$ sea fibrada, es necesario y suficiente que los
$c_{f,g}$ sean isomorfismos.
\end{proposition}

Tomando para~$f$ un isomorfismo y para~$g$ su inverso, y considerando los
isomorfismos $c_{f,g}$ y $c_{g,f}$, se concluye:
\begin{corollaire}
\label{VI.7.3} Si $\cal{F}$ es una categoría \emph{fibrada con clivaje}
sobre~$\cal{E}$, entonces, para todo isomorfismo $f\colon T\to S$ en
$\cal{E}$, $f^*$ es una equivalencia de categorías
$\cal{F}_S\to\cal{F}_T$.
\end{corollaire}

\begin{proposition}
\label{VI.7.4} Sea $\cal{F}$ una categoría con clivaje sobre~$\cal{E}$.
Se tiene

$\qquad\qquad A)\qquad
\left\{\begin{array}{ccc}c_{f,\id_T}(\xi)&=&\alpha_{\id_T}(f^*(\xi))\\
c_{\id_S,f}(\xi)&=&f^*(\alpha_{\id_S}(\xi))
\end{array}\right.$

\smallskip
$\qquad\qquad B)\qquad c_{f,gh}(\xi)\cdot
c_{g,h}(f^*(\xi))=c_{fg,h}(\xi)\cdot h^*(c_{f,g}(\xi))$

\noindent(En
\marginpar{173}
estas fórmulas, $f,g,h$ denotan morfismos
$$
V\to U\to T\to S
$$
y $\xi$ un objeto de~$\cal{F}_S$).
\end{proposition}

La primera y la segunda relación, en el caso de un clivaje normalizado,
adoptan la forma más sencilla
\begin{equation*}
\label{eq:VI.7.A'}\quad A')\qquad {c_{f,\id_T}=\id_{f^*},\quad
c_{\id_S,f}=\id_{f^*}.\qquad\qquad\qquad\qquad\qquad\qquad}
\end{equation*}

En cuanto a la tercera, puede visualizarse mediante la conmutatividad del
diagrama
\begin{equation*}
\label{eq:VI.7.D} \tag{$D$}
\begin{split}
\xymatrix@C=2.2cm{
h^*g^*f^*(\xi) \ar[r]^-{c_{g,h}(f^*(\xi))} \ar[d]_{h^*(c_{f,g}(\xi))}
& (gh)^*(f^*(\xi)) \ar[d]^{c_{f,gh}(\xi)} \\
h^*(fg)^*(\xi) \ar[r]_-{c_{fg,h}(\xi)} & (fgh)^*(\xi) }
\end{split}
\end{equation*}
En el caso de las categorías
\ifthenelse{\boolean{orig}}
{fibradas,}
{fibradas}
(donde los $c_{f,g}$ son isomorfismos), esta conmutatividad puede expresarse
intuitivamente por el hecho de que \emph{la utilización sucesiva de los
isomorfismos de la forma $c_{f,g}$ no conduce a \og identificaciones
contradictorias\fg}. También puede escribirse esta fórmula sin el argumento
$\xi$, utilizando el producto de convolución de transformaciones naturales:
$$
c_{fg,h}\circ (h^**c_{f,g})=c_{f,gh}\circ(c_{g,h}*f^*).
$$

\ifthenelse{\boolean{orig}}
{\ignorespaces}
{\enlargethispage{.5cm}}
La demostración de las dos primeras fórmulas de~\Ref{VI.7.4} es trivial;
esbocemos la de la tercera. Para ello, consideremos, además del cuadrado
$(D)$, el cuadrado de transformaciones naturales:
\begin{equation*}
\label{eq:VI.7.D'} \tag{$D'$}
\begin{split}
\xymatrix@C=2cm{
g^*f^*(\xi) \ar[r]^-{\alpha_g(f^*(\xi))}
\ar[d]_{c_{f,g}(\xi)} & f^*(\xi) \ar[d]^{\alpha_f(\xi)} \\
(fg)^*(\xi) \ar[r]_-{\alpha_{fg}(\xi)} & \xi }
\end{split}
\end{equation*}
que
\marginpar{174}
es conmutativo por definición de $c_{f,g}(\xi)$. Consideremos el diagrama
obtenido uniendo los vértices de~$(D)$ a los vértices correspondientes de
$(D')$ mediante las transformaciones naturales de la forma~$\alpha$:
$$
\begin{array}{ll}
\alpha_h(g^*f^*(\xi)),&\alpha_{gh}(f^*(\xi))\\
\alpha_h((fg)^*(\xi)),&\alpha_{fgh}(\xi).
\end{array}
$$
Las cuatro caras laterales del cubo así obtenido son también conmutativas:
para la de la izquierda, esto se debe a que la columna izquierda de~$(D)$
se deduce de la columna izquierda de~$(D')$
\ifthenelse{\boolean{orig}}{por aplicación de~$h$}{por aplicación de~$h^*$}%
, y a que $\alpha_h$ es una transformación natural; para las otras tres,
no es sino la definición de las operaciones~$c$ de los tres lados restantes
de~$(D)$. Así, las cinco caras del cubo distintas de la cara superior son
conmutativas. Se sigue que los dos $(fgh)$-morfismos
$h^*g^*f^*(\xi)\to(fgh)^*(\xi)$ definidos por~$(D)$ tienen el mismo
compuesto con $\alpha_{fgh}(\xi)\colon(fgh)^*(\xi)\to\xi$; por tanto, son
iguales por definición de~$(fgh)^*$.

Para lo que sigue, limitémonos a las categorías con clivaje
\emph{normalizado}. Una categoría tal da lugar a los objetos siguientes:
\begin{enumerate}
\item[a)] Una aplicación $S\mto\cal{F}_S$ de~$\Ob(\cal{E})$ en~$\Cat$.
\item[b)] Una aplicación $f\mto f^*$ que asigna a todo
$f\in\Fl(\cal{E})$, de fuente~$T$ y destino~$S$,
\ifthenelse{\boolean{orig}}
{una}
{un}
funtor $f^*\colon\cal{F}_S\to\cal{F}_T$.
\item[c)] Una aplicación $(f,g)\mto c_{f,g}$ que asigna a todo par de
flechas $(f,g)$ de~$\cal{E}$ una transformación natural
$c_{f,g}\colon g^*f^*\to(fg)^*$.
\end{enumerate}

Además, estos datos satisfacen las condiciones expresadas en las fórmulas
$A')$ y~$B)$ dadas más arriba. (N.B. Si no nos hubiéramos limitado al caso
de un clivaje normalizado, habría sido necesario introducir un dato
suplementario, a saber, una función $S\mto\alpha_S$ que asocia a todo objeto
$S$ de~$\cal{E}$ una transformación natural
$\alpha_S\colon(\id_S)^*\to\id_{\cal{F}_S}$; la condición $A')$ sería
entonces sustituida por la condición~$A)$).

Vamos
\marginpar{175}
a mostrar ahora cómo puede reconstruirse (salvo isomorfismo único) la
categoría con clivaje normalizado $\cal{F}$ sobre~$\cal{E}$ a partir de los
objetos precedentes.
